\section*{Open Questions}

The following five questions arose from doing this replication (not paraphrased from
Watrous's own ``future directions'' -- he wrote his in 2000 and most of those are now
resolved). Each is grounded in a specific observation from our simulation.

\begin{enumerate}
\item[Q1.] \textbf{Empirical shot budget for HSP recovery is essentially flat ($\le 4$) across
$|G|\in[6,15]$ in our runs, despite $t$ growing.} The theoretical bound is polylog $|G|$
but has non-trivial constants inside the continued-fraction success probability. Where is
the crossover to the theoretical shot growth? A systematic sweep $N\in[16, 2^{20}]$ with
Aer statevector on 20+ qubits would pin the empirical polylog exponent.

\item[Q2.] \textbf{Watrous's oracle assumes exact group operations at unit cost, but
noisy near-term hardware would introduce coset-drift errors that compound through the
recursion.} How does the trace distance of $|G\rangle$ scale when the black-box oracle
itself is depolarized at rate $p$ per query, as a function of derived-series length and
$|G|$? This is now testable on real superconducting hardware for $|G|\le 32$ and would
give a concrete resource estimate for the algorithm's near-term applicability.

\item[Q3.] \textbf{Our $D_4$ coset-decomposition is exact only because the encoding is
lossless (3 qubits, 8 basis states, permutation encoding).} For a black-box group presented
by permutations of $\{0,\ldots,m-1\}$ with $m > |G|$ (Watrous's actual setting), preparing
$|G\rangle$ requires projecting out non-group basis states. What is the minimum quantum
resource cost (ancillas + queries) of this projection when the group is presented via a
random Cayley-table oracle vs.\ a structured (polycyclic) presentation?

\item[Q4.] \textbf{The QFT-based HSP subroutine we implemented uses $t = 2\log_2 N$
qubits, but Watrous never gives a tight bound on the required $t$ for the recursive
composition.} Empirically we saw $t = 2\log_2 N$ suffices at every level, but does
error compounding across derived-series depth $\ell$ force $t = O(\ell \log|G|)$?
A benchmark on solvable groups of controlled derived length (e.g., iterated wreath
products $\mathbb{Z}_2 \wr \mathbb{Z}_2 \wr \cdots$) would resolve this.

\item[Q5.] \textbf{The HSP oracle we built is a permutation \texttt{UnitaryGate} of size
$2^{t+\lceil\log_2 d\rceil} \times 2^{t+\lceil\log_2 d\rceil}$}, which is exactly the
``black-box'' cost model Watrous assumes but is $\Omega(|G|^2)$ classical memory.
For solvable groups given by a polycyclic presentation (typical in computational algebra
systems), a much more compact reversible-circuit compilation of the oracle should exist.
What is the smallest T-count / gate-count reversible circuit for the group operation of
the smallest non-abelian solvable groups ($D_n$ for $n\le 8$, $Q_8$, $S_3$, $A_4$), and
does it beat the $\Theta(|G|^2)$ naive lookup?
\end{enumerate}
